\subsubsection*{Meses extremos y contraste con las series cronológicas}
Se identifican los mínimos y máximos de los mismos 285 meses completos de 1998--2022 y se contrastan con las series de la Figura 8 y los valores de sus meses vecinos. Son extremos de la muestra común, que se distingue de los registros completos por fuente. Los meses adicionales con valores altos se evalúan mediante P95 y los límites de las cajas; no se exige que todas las variables sean extremas simultáneamente.

\begin{center}
{\small
\begin{tabular}{lll}
\toprule
Variable & Mínimo (mes; valor) & Máximo (mes; valor) \\
\midrule
CHIRPS (mm/mes) & 2001-08; 30,12 & 2022-10; 407,57 \\
IMERG polígono (mm/mes) & 2015-12; 60,83 & 2010-11; 440,52 \\
Caudal Q (m$^3$/s) & 2019-09; 20,33 & 2010-11; 508,77 \\
Escorrentía R (mm/mes) & 1998-02; 18,45 & 2010-11; 471,45 \\
ERA5-Land Tmedia ($^\circ$C) & 2000-01; 15,10 & 1998-01; 17,72 \\
\bottomrule
\end{tabular}
}
\end{center}

\textbf{Noviembre de 2010.} Coinciden los máximos de IMERG, Q y R; CHIRPS también registra lluvia muy alta (396,03 mm/mes). Q pasa de 216,17 m$^3$/s en octubre a 508,77 en noviembre y permanece en 439,72 en diciembre. Ambas precipitaciones son elevadas en esos meses. La secuencia es compatible con un episodio húmedo sostenido y una respuesta prolongada de la cuenca, más que con un pico aislado; no constituye una validación independiente del dato.

\textbf{Noviembre de 1999.} CHIRPS e IMERG registran 277,91 y 319,05 mm/mes, respectivamente. Q alcanza 328,46 m$^3$/s y R, 304,36 mm/mes: ambos superan P95 y los límites superiores de las cajas. Octubre ya presenta lluvia elevada (259,15 y 314,74 mm/mes) y diciembre conserva Q de 261,25 m$^3$/s. Es una secuencia plausible de acumulación de humedad y caudal alto, aunque las precipitaciones de noviembre no superan sus respectivos P95.

\textbf{Octubre de 2022.} El máximo CHIRPS coincide con IMERG elevado (358,75 mm/mes). Q aumenta de 73,93 m$^3$/s en septiembre a 204,28 en octubre y 296,50 en noviembre; ambas precipitaciones permanecen altas en noviembre. Es compatible con almacenamiento y respuesta retardada; la comparación mensual no determina el tiempo de respuesta exacto.

\textbf{Meses con valores bajos y extremos de temperatura.} En agosto de 2001, el mínimo CHIRPS coincide con IMERG de 96,98 mm/mes y Q bajo (29,21 m$^3$/s). En enero de 1998, la temperatura máxima coincide con poca lluvia (50,72 mm/mes en CHIRPS y 75,91 en IMERG) y Q de 22,29 m$^3$/s; febrero mantiene Q bajo (21,34 m$^3$/s) y registra el mínimo R. En enero de 2000, la temperatura mínima (15,10 $^\circ$C) sigue a los meses lluviosos de finales de 1999, con Q todavía elevado (199,26 m$^3$/s); la coincidencia es plausible, pero no demuestra una causa climática. En septiembre de 2019, Q mínimo sucede después de agosto, con poca lluvia (40,56 mm/mes en CHIRPS y 89,28 en IMERG). La lluvia se recupera en septiembre, mientras Q aumenta en octubre a 54,01 m$^3$/s: no se espera necesariamente una respuesta simultánea.

\textbf{Diciembre de 2015: discrepancia que requiere revisión.} IMERG alcanza su mínimo (60,83 mm/mes), mientras CHIRPS registra 170,66 mm/mes. Q es bajo (29,70 m$^3$/s) y la temperatura alta (17,28 $^\circ$C); enero de 2016 mantiene poca lluvia, Q bajo y temperatura alta. La situación hidrológica es plausible, pero la diferencia entre productos impide considerar validado el mínimo IMERG solo por su coincidencia con Q. Debe contrastarse con información independiente antes de atribuirla a un error o a un episodio regional concreto.

\textbf{Conclusión y alcance.} Los extremos presentan coherencia general con sus meses vecinos; esta revisión no demuestra errores ni certifica exactitud. No se elimina ningún dato. Q y R no aportan evidencias independientes, porque R deriva de Q; la duración del mes explica que sus mínimos no coincidan necesariamente. La temperatura procede de reanálisis y tampoco sustituye una observación local. En el registro completo de IMERG, el máximo es noviembre de 2008 (447,12 mm/mes), excluido de la muestra común porque faltan CHIRPS y Q: su coherencia con esas variables no puede evaluarse. Los extremos históricos anteriores a 1998, ya identificados en las tablas del registro completo, requieren contraste adicional con sus cronologías; este análisis se limita al periodo común y no atribuye los episodios a eventos climáticos sin evidencia externa.

% EXTREMOS_HISTORICOS_INICIO
\subsubsection*{Extremos del registro completo anteriores a 1998}
\begin{center}{\small
\begin{tabular}{lllr}
\toprule
Variable & Unidad & Mes mínimo & Mínimo \\
\midrule
CHIRPS & mm/mes & 1982-08 & 24.89 \\
Caudal Q & m³/s & 1992-07 & 18.37 \\
Escorrentía R & mm/mes & 1992-07 & 17.59 \\
ERA5-Land Tmedia & °C & 1984-01 & 14.84 \\
\bottomrule
\end{tabular}
}\end{center}
Este contraste utiliza el registro completo por fuente, separado de los 285 meses comunes con IMERG. Se revisan agosto de 1982, julio de 1992 y enero de 1984 porque contienen los mínimos históricos que no aparecen en la figura de 1998--2022. Son datos válidos anteriores a la cobertura IMERG utilizada, no meses descartados por ese motivo. Los máximos de CHIRPS, Q, R y temperatura están dentro del periodo ya analizado; el máximo IMERG de noviembre de 2008 se mantiene aparte por ausencia de CHIRPS y Q.

Agosto de 1982: CHIRPS alcanza 24,89 mm/mes, tras 87,29 en junio y 53,49 en julio. Q desciende de 77,22 a 45,62 y 29,20 m$^3$/s; R es 27,96 mm/mes en agosto. En septiembre aumenta la lluvia a 168,27 mm/mes, pero Q permanece bajo (31,70 m$^3$/s), y en octubre sube a 66,39 m$^3$/s. Es una secuencia plausible de condiciones secas y recuperación posterior; el mínimo de lluvia no constituye por sí solo un problema de datos.

Julio de 1992: Q y R alcanzan sus mínimos históricos (18,37 m$^3$/s y 17,59 mm/mes). CHIRPS registra 75,46 mm/mes, después de 80,03 en junio; Q ya era bajo en mayo y junio (28,16 y 28,38 m$^3$/s). La persistencia de caudal bajo es compatible con almacenamiento reducido, pero agosto y septiembre tienen datos ausentes de CHIRPS, Q y R: no puede comprobarse la evolución posterior con esas series ni rellenarse el vacío con cero.

Enero de 1984: ERA5-Land alcanza su temperatura mínima (14,84 $^\circ$C), con valores próximos en diciembre (15,12 $^\circ$C) y febrero (15,05 $^\circ$C). CHIRPS permanece entre 190,87 y 236,26 mm/mes en esos tres meses y Q entre 98,77 y 127,13 m$^3$/s. El mínimo está integrado en una secuencia de temperaturas bajas y lluvia apreciable, sin un salto aislado evidente. Es plausible, pero no prueba causalidad ni exactitud del reanálisis. En conjunto, no se identifican errores confirmados: se conservan todos los valores y se reconoce la limitación de los faltantes de 1992. R deriva de Q, por lo que su coincidencia no constituye evidencia independiente.
\begin{figure}[H]\centering\includegraphics[width=\linewidth]{figuras/extremos_historicos.pdf}\caption{Extremos históricos y meses vecinos. La línea discontinua señala el mes extremo; los cortes conservan las ausencias. R se detalla en el texto y la tabla.}\end{figure}
% EXTREMOS_HISTORICOS_FIN

% CONCLUSION_EXTREMOS_INICIO
\subsubsection*{Conclusión del contraste de extremos}
\textbf{¿Son coherentes entre variables?} Sí, en general: los meses lluviosos coinciden con caudales altos o con aumentos posteriores, y los meses secos con caudales bajos. La principal discrepancia aparece en diciembre de 2015, cuando IMERG registra mucha menos lluvia que CHIRPS.

\textbf{¿Son episodios plausibles o problemas de datos?} Las secuencias son físicamente plausibles y no encontramos evidencia suficiente para declarar errores. La discrepancia de diciembre de 2015 requiere revisión independiente; los faltantes posteriores a julio de 1992 limitan su comprobación.
% CONCLUSION_EXTREMOS_FIN
